% ============================================================
% Chapter 7: Results
%   Section: Pilot on X-Fact
%     Subsection: What the Years Between Claim and Search Do
% Included from chapters/07-results/01-xfact-pilot/xfact-pilot.tex
% ============================================================

\subsection{What the Years Between Claim and Search Do}
\label{subsec:pilot-time}

The English claims were made between 2008 and 2018, and X-Fact recorded
no date for the Spanish ones; the dataset was published in 2021; the web
was searched in October 2026; and the model's pretraining data ends in
December 2023, according to its model card. A label that was right
when a fact-checker wrote it is being compared with what the web says
years later. Six effects follow. Each is counted on the pilot, with the
accuracy of the claims it touches against the rest; at this size those
differences are indications, not findings.

\paragraph{Verdict leakage.} The fact-check itself is now on the web.
PolitiFact's own verdict page was ranked for 9 of the 24 English claims
and cited for 8, four of them from PolitiFact's staging server, which
search engines index. Chequeado's page was never retrieved for any of
the 40 Spanish claims, although every row links it and all 42 of those
pages still answer: with Google refusing the metasearch engine, the
engines that answered did not surface Chequeado. The leak inflated less
than feared, because the small model read the verdict page as evidence
rather than as a ruling: 4 of the 9 leaked claims were right (44\%,
against 27\% for the rest). A stronger model would read the ruling.

\paragraph{Temporal leakage.} Of the 20 claims for which it can be
decided, 16 cited a source published after the claim. The custom set's
annotation guide forbids exactly this (Rule~3), but the pipeline
searches without a date limit, and on X-Fact newer evidence is the norm:
leaving out the claims with either leak lowers the pilot's accuracy
from 29.7\% to 23.3\%.

\paragraph{Drift.} Twenty of the 64 claims tie themselves to their own
present, twelve with relative time (``hoy'', ``now'', ``este año'') and
nine with records or rankings (``de la historia'', ``the first time in
history''). Without a claim date, ``hoy'' means the day of the search: a
claim that the Asignación Universal por Hijo ``today'' reaches four
million beneficiaries, true when Chequeado checked it, was judged false
on articles from 2026; a 2012 claim about the worst cities for asthma
was judged true on the 2024 and 2026 rankings. No accuracy penalty
shows at this size, and none can be measured properly without a
re-dated, re-labelled copy of each claim.

\paragraph{Link rot.} Of the pilot's 305 reference links, the 120
English ones were never available (X-Fact ships PolitiFact's rows
without them). Of the 69 Spanish links that are neither Chequeado's
nor social-media posts, 14 no longer exist (404, a domain that no
longer resolves, or a redirect to a home page) and 4 refused the check;
5 of the 14 survive in the Wayback Machine. The system retrieved none
of the annotator's links, so recall measured in 2026 is bounded by what
has survived, and by search engines that no longer reach the web of
2017.

\paragraph{Recency of search.} Of the 235 ranked sources with a
publication date, 53\% appeared after X-Fact was built and 23\% in 2026,
the year of the search, the most frequent year of all. For the dated
English claims the median ranked source is 6.8 years newer than the
claim. The evidence a fact-checker read at the time is a minority of
what search returns now.

\paragraph{The model's memory.} Every claim and every fact-check of it
predates the model's training cutoff, so the verdicts could be in its
weights. Asked about each claim with no evidence, the model reported
remembering a fact-check for none of the 64, and answered
\texttt{FALSE} to 60 of them, scoring below the majority baseline.
There is no sign that this 3-billion-parameter model recalls X-Fact's
verdicts; a larger or more recent model should be probed the same way
before its X-Fact score is read.

\paragraph{Consequence.} For benchmark runs on fact-checker datasets,
the fact-checkers' own domains should be excluded at retrieval time,
through a per-run threshold that is off by default; in production, and
on the custom set, a published fact-check is legitimate evidence and
must stay available. This is recorded as a recommendation and has not
been implemented.
